\subsection{应用：III. 广义幂}

设 E 为具有单位元 e 且合成律记为 T 的幺半群。若 x 在 E 中可逆，令 \(g_{x}\) 为将 1 映到 x 的 Z 到 E 的同态。令 \(g_{x}(n) = \frac{n}{T} x\) （对所有 \(n \in Z\) ）；由引理2，此记号在 \(n \in N\) 时与先前引入的记号相容。则

\[
{ } ^ { m + n } \top x = \left( \begin{array} { c } m \\ T \end{array} x \right) \top \left( \begin{array} { c } n \\ T \end{array} x \right)\tag{9}
\]

\[
\frac {0}{T} x = e\tag{10}
\]

\[
\frac {1}{T} x = x\tag{11}
\]

对于可逆的 \(x\) （在 \(\mathbf{E}\) 中）和 \(m, n\) （在 \(\mathbf{Z}\) 中）。进一步，若 \(y = \frac{m}{T} x\) ，则映射 \(n \mapsto g_x(mn)\) （从 \(\mathbf{Z}\) 到 \(\mathbf{E}\) ）是将 1 映到 \(y\) 的一个同态，因此 \(g_x(mn) = g_y(n)\) ，即

\[
\frac {m n}{T} x = \frac {n}{T} \left(\frac {m}{T} x\right).\tag{12}
\]

由于 \(-1\) 是 \(1\) 在 \(\mathbf{Z}\) 中的负元， \(\overline{\top}^{-1}x\) 是 \(x = \frac{1}{\top} x\) 在 E 中的逆元。若在(9)中记 \(n = -m\) ，则可见 \(\overline{\top}^{-m}x\) 是 \(\overline{\top}^{m}x\) 的逆元.

